\begin{textsample}{POS Dim 6 – global\_north\_2023\_11 – Score 4.00 – global\_north\_2023\_11/103898187.txt}  \label{ex:f6_pos_218}
Hospital and council staff in vigil to ' highlight health crisis ' in Gaza Health workers at DCH hold vigil for victims of Gaza Health care and council workers came together to take part in a vigil to support the people of Gaza.
The group came together today, Wednesday, November 29, in conjunction with the United Nations International Day of Solidarity with the Palestinian people.
The group, consisting of doctors, nurses and other medical professionals as well as social workers, were staging the protest in solidarity with Gaza as 25 out of 36 hospitals in the territory have ceased functioning in Israel ' s war with Hamas.
They gathered outside Dorset County Hospital in Dorchester from 12pm with placards, flags and signs in protest of the war.
Phillip Marfleet, of the Dorset Palestinian Solidarity Campaign, said :"We ' re at the hospital to highlight the crisis the health system in Gaza because the great majority of the hospitals have been destroyed or rendered inoperative.
Very large numbers of doctors and other heath staff have been killed and injured @ @ @ @ @ @ @ @ @ @ at our own hospital want to have an act of solidarity and call for a permanent ceasefire which means that this sort of horrible violence comes to an end."Lynne Hubbard, a dietician at DCH, said :"As health workers, we ' re very aware of what ' s happening in Gaza, health workers have stayed with their patients."13,000 Palestinians are dead, 70 per cent of those are women and children and there is nowhere for people of Gaza to go."When we saw the tragedy in Ukraine, people could go to Poland and to Europe, but it ' s the size of the Isle of Wight, and 15 times the population.
As humans and as health workers we feel like we have to raise a voice.""We ' ve been round wards talking to people and they ' re very supportive.
Health workers are all humanitarians, and we all save lives.
We keep thinking would we stay and risk our lives."She went on @ @ @ @ @ @ @ @ @ @ informed as the workers made it clear it would be a peaceful event and there was no objection to the protest.
Dr Muhammed Bayoumy, a consultant at DCH, addressed the protestors saying :"When this atrocity first happened, I was feeling hopeless and my spirit was down.
I was afraid that staying here would mean I watch what happened in Palestine without any action of solidarity."These many protests have been keeping me going, knowing there are many people in Dorchester with kind hearts keeps me going."Students from Thomas Hardye Sixth Form also marched down on their lunch break to join the vigil where they were congratulated by Mr Marfleet and the protestors for their march last week.
The protest lasted around two hours and \textbf{campaigners} have planned another vigil for Saturday, December 3 at the town pump which they hope will see hundreds of protesters join them.
Comments : Our rules We want our comments to be a lively and valuable part of our community - a place @ @ @ @ @ @ @ @ @ @ local issues.
The ability to comment on our stories is a privilege, not a right, however, and that privilege may be withdrawn if it is abused or misused. follow us This website and associated newspapers adhere to the Independent Press Standards Organisation ' s Editors ' Code of Practice.
If you have a complaint about the editorial content which relates to \textbf{inaccuracy} or \textbf{intrusion}, then please contact the editor here.
If you are \textbf{dissatisfied} with the response provided you can contact IPSO here

% matched lemmas: campaigner, dissatisfy, inaccuracy, intrusion
\end{textsample}
